% Generated draft: requires mathematical and source review.
\documentclass[11pt]{article}
\usepackage[margin=25mm]{geometry}
\usepackage{fontspec}
\setmainfont{Latin Modern Roman}
\usepackage{amsmath,amssymb,mathtools,unicode-math}
\setmathfont{Latin Modern Math}
\usepackage{enumitem,needspace,xurl,hyperref}
\setlength{\parindent}{0pt}
\setlength{\parskip}{5pt}
\setlength{\emergencystretch}{3em}
\newcommand{\N}{\mathbb{N}}
\newcommand{\Z}{\mathbb{Z}}
\newcommand{\Q}{\mathbb{Q}}
\newcommand{\R}{\mathbb{R}}
\newcommand{\C}{\mathbb{C}}
\newcommand{\F}{\mathbb{F}}
\newcommand{\A}{\mathbb{A}}
\newcommand{\PP}{\mathbb{P}}
\newcommand{\E}{\mathbb{E}}
\newcommand{\eps}{\varepsilon}
\newcommand{\dd}{\,\mathrm{d}}
\begin{document}
\section*{UnsolvedMath ID: 30001494}
\textbf{Source problem code:} catalogue entry
\subsection{Definitions and mathematical statement}
\textbf{Supplied question.} If \(i\colon Y\hookrightarrow X\) is a smooth subvariety whose normal bundle \(N_{Y/X}\) is ample, must its cycle class \([Y]\) be big?

For a regular closed immersion with ideal sheaf \(\mathcal I_Y\), the normal bundle is
\[
N_{Y/X}=\mathcal{H}om_{\mathcal O_Y}(\mathcal I_Y/\mathcal I_Y^2,\mathcal O_Y).
\]
A vector bundle \(E\) on \(Y\) is ample when the tautological quotient line bundle \(\mathcal O_{\mathbb P_Y(E)}(1)\) is ample, using the convention that \(\mathbb P_Y(E)\) parametrizes one-dimensional quotients.

A standard numerical interpretation of the requested conclusion can be specified conditionally as follows. If \(X\) is a projective variety over a field \(k\) and \(Y\) is a closed subvariety of pure dimension \(d\), let \(N_d(X)_{\mathbb R}\) be the real vector space of \(d\)-dimensional algebraic cycles modulo numerical equivalence. Let
\[
\overline{\mathrm{Eff}}_d(X)
=\overline{\left\{\sum_{j=1}^{m}a_j[Z_j]\;:\;m\geq 0,\ a_j\in\mathbb R_{\geq 0},\ Z_j\subseteq X\text{ integral},\ \dim Z_j=d\right\}}
\subseteq N_d(X)_{\mathbb R}.
\]
In this interpretation, the question asks whether
\[
N_{Y/X}\text{ ample}
\quad\Longrightarrow\quad
[Y]\in\operatorname{int}_{N_d(X)_{\mathbb R}}\bigl(\overline{\mathrm{Eff}}_d(X)\bigr).
\]
Here the interior is taken in the usual topology of the real numerical cycle space.

\textbf{Unresolved qualifications.} The supplied statement does not specify the base field, the ambient hypotheses on \(X\), whether the embedding is closed and regular, the dimension or codimension restrictions, or the equivalence relation intended for cycle classes. In particular, projectivity and smoothness of \(X\) are not stated. The numerical formulation above explains a conventional meaning under explicit additional hypotheses; it does not establish that these are the intended source hypotheses. The record also does not settle whether zero-dimensional subvarieties or \(Y=X\) are included. Consequently, the universally quantified domain of the question requires source clarification.
\subsection{Short English statement}
The question asks whether a smooth subvariety with ample normal bundle necessarily has a big cycle class. Under the usual projective numerical interpretation, this means that its class lies in the interior of the closed cone generated by effective cycles of the same dimension. The supplied record leaves the ambient hypotheses and the intended cycle-class interpretation unresolved.
\subsection{Sources}
\begin{itemize}
\item Supplied local UnsolvedMath registry record, permanent ID \(30001494\), title \emph{Bigness of Subvarieties with Ample Normal Bundle}, exact source problem code \texttt{catalogue entry}, in the Oberwolfach Reports Open Problems collection; proposed year recorded as 2010.
\item \emph{Classical Algebraic Geometry} (2010), Oberwolfach Reports, DOI 10.4171/owr/2010/27, pp.\ 1573--1623. Supplied source citation and literature-source label: ``original OWR Report 27/2010.'' \url{https://doi.org/10.4171/owr/2010/27}. No problem-specific page is supplied.
\item Supplied literature source: \emph{Subvarieties with ample normal bundle}. \url{https://ems.press/content/serial-article-files/32151?nt=1}. The record supplies no authors or further bibliographic details for this link.
\item The supplied assessment calls this Peternell's question, although the registry's proposer field is empty. The supplied literature status is \emph{partially solved}. Its assessment states that affirmative results exist for curves and several special settings, while the question remains open for arbitrary smooth subvarieties with ample normal bundle. Literature checked at: \texttt{2026-08-21}. This is supplied dated curation, not an independently verified literature assessment; the record explicitly requests independent verification of time-sensitive status and cited claims.
\end{itemize}
\end{document}
